% =====================================================================
% Appendix B -- Numerical methods (notes 4.8 and 4.9 point 5). Written in
% step 6.4. Implementation: beltloop/src/beltloop (eigen.py, transfer.py,
% forcing.py, response.py, lumped.py); numbers in paper_numbers.tex
% (section "Appendix B").
% =====================================================================
\section{Numerical methods}\label{app:numerics}

The closed-form solution is implemented in the open-source Python package
\texttt{beltloop} (see Data availability), together with the lumped-mass model of
Section~\ref{sec:lumped} and the scripts that produce every figure, table and number
of the paper. This appendix describes the numerical steps that the closed form leaves
open, and the lumped-mass model used for its verification.

\paragraph{Natural frequencies.} The roots of~\eqref{eq:chareq} are found without a
search grid, from the interlacing property of Section~\ref{sec:properties}. The poles
of~\eqref{eq:receptance}, the fixed--free frequencies of each strand, are bracketed
first with the scaled Pr\"ufer angle $\phi=\operatorname{atan2}(\kappa W,W')$ of the
chain started from $(W,W')=(0,1)$ at the drive. Inside a segment $\phi$ grows by
$\kappa\ell$; at a change of density $W$ and $W'$ are continuous, so $\tan\phi$ is
multiplied by the ratio of the wavenumbers and $\phi$ is mapped within its quadrant.
The end value $\phi(\Omega)$ increases with $\Omega$ (Sturm), and the $j$-th
fixed--free frequency is the root of $\phi(\Omega)=(j-\tfrac12)\pi$. Since each of the
$m$ density changes of a strand moves $\phi$ by less than $\pi/2$, that root lies in
\begin{equation}\label{eq:prufer_bracket}
  \frac{(j-\tfrac12)\pi-\tfrac{m}{2}\pi}{\Phi}\le\Omega_j\le\frac{(j-\tfrac12)\pi+\tfrac{m}{2}\pi}{\Phi},
  \qquad \Phi=\sum_i\kappa_i\ell_i/\Omega,
\end{equation}
which is used as the bracket for Brent's method. The poles of the two strands are then
merged and sorted, $p_1\le p_2\le\dots$, and the $k$-th natural frequency is sought in
$(p_{k-1},p_k)$, $p_0=0$. In that interval $D/(\beta\Omega^2A_{22}B_{11})$ increases
from $-\infty$ to $+\infty$ and $A_{22}B_{11}$ keeps its sign, so $D$ multiplied by that
sign changes from negative to positive across the interval; the end values are replaced
by these limit signs, and Brent's method converges to the single root. A pair of
coincident poles is itself a root and is returned exactly. The method finds every
root, including pairs separated by $O(\beta)$ near nearly coincident poles, which a grid
search with a fixed step misses when $\beta$ is small. The limit $\beta=0$ returns the
poles themselves. For the drive without speed control, Eq.~\eqref{eq:chareq_torque},
the interlacing argument is not available: real roots are located by a sign scan on a
fine grid, and the complex roots with a slip dashpot by Newton's method started from
the root with prescribed velocity, the limit $\hat c\to\infty$. Both serve only for the
checks of Section~\ref{sec:fieldchecks}.
% tests: test_eigen.py (fixed-free roots, all roots against a dense scan, close roots
%        at beta = 1e-7, FE agreement, torque variant), test_harrison.py (damped drive)

\paragraph{Modes and modal quantities.} Each mode is obtained by propagating
$(W,W')=(0,1)$ through the chain with the matrices~\eqref{eq:transfer}, with the jump of
$W$ and the take-up amplitude $Y$ of~\eqref{eq:eigen_tu} at the take-up. On a segment
$W$ is a combination of $\cos\kappa x$ and $\sin\kappa x$, so the integrals that enter
the modal mass~\eqref{eq:inner}, the participation factors~\eqref{eq:participation}
and the stiffness identity $\int W'^2\dd x=\Omega^2m$ are evaluated in closed form,
segment by segment; the last identity is kept as a check of every mode. The modes are
normalised to unit modal mass. The take-up term $\beta Y^2$ of the modal mass requires
$\beta>0$; the limit $\beta\to0$ of the forced response is computed with a small
$\beta$, which the free-end structure makes harmless.
% tests: test_modes.py (interfaces, closed-form integrals, weighted orthogonality,
%        Parseval, quasi-static sums)

\paragraph{Exact integration of the modal equations.} On each interval of a start-up
profile the loads $\hat a(\tau)$ and $\varphi(\tau)$ are outputs of a small linear
system $\mat z'=\mat E\mat z$: the polynomial pieces of the triangular, parabolic and
piecewise-linear profiles, and the sine and cosine of the sinusoidal profile, whose
velocity gives the onset $\varphi=V/V_\infty$. Appending $\mat z$ to the state
$(p_k,p_k')$ of each mode gives an autonomous linear system, which is advanced exactly
to every output time by its matrix exponential, interval by interval, with the modal
state carried across the breaks. The integration has no time step and no truncation
error; it holds for any damping, including overdamped high modes, and for exact
resonance $\Omega_k=\varpi$. A jump of the drive velocity, as in the offset start of
\citet{lodewijks1996}, is replaced by a short ramp so that the load stays bounded.
% tests: test_response.py (exosystem against closed form, exact integration against an
%        ODE solver, residual amplitude at exact resonance, damping ratio)

\paragraph{Quasi-static split.} The modal sum~\eqref{eq:outputs} for the tension
converges slowly, because the tension is the derivative of the displacement and the
loads reach every mode. It is therefore split into the quasi-static
field~\eqref{eq:quasistatic}, summed in closed form, and a dynamic remainder. With
Kelvin--Voigt damping the quasi-static tension is exact but the strain lags; the
quasi-static modal coordinates are accordingly
$p_k^{\mathrm{qs}}=-(\Gamma_k\hat a_f+R_k\varphi_f)/\Omega_k^2$, where the lagged loads
obey $2\hat\zeta\hat a_f'+\hat a_f=\hat a$ and $2\hat\zeta\varphi_f'+\varphi_f=\varphi$,
a first-order filter with the same time constant for every mode, and
\begin{equation}\label{eq:split}
  \hat T=\hat T_{\mathrm{qs}}+\sum_k\left(q_k+2\hat\zeta q_k'\right)W_k',\qquad
  \hat y=\tfrac12\bigl(\hat a_fI_\mu+\varphi_fI_r\bigr)+\sum_kq_kY_k,\qquad
  q_k=p_k-p_k^{\mathrm{qs}},
\end{equation}
with $I_\mu=\int_0^2\int_\xi^x\hat\mu\,\dd x'\,\dd x$ and $I_r$ defined likewise. The
remainder $q_k$ is driven only by the rate of change of the quasi-static response and
converges fast with or without damping. With \NumSplitModes\ modes the tension differs
from a reference with \NumSplitRefModes\ modes by at most \num{\NumSplitErr} of its
maximum for $0\le\hat\zeta\le0.1$, against \num{\NumPlainErr} for the plain sum with the
same modes. Splitting off the undamped quasi-static response instead, $p_k-F_k/\Omega_k^2$,
leaves a remainder of order $2\hat\zeta F_k'/\Omega_k^2$, which converges as slowly as
the plain sum.
% tests: test_response.py (split convergence with damping, lagged loads, quasi-static
%        limit); numbers: paper_numbers.py, section "Appendix B"

\paragraph{Lumped-mass model.} The model of Section~\ref{sec:lumped} divides the loop,
from the drive exit to the drive entry, into $N$ elements with nodes at the head and
tail pulleys and at the take-up. Each element carries half of its mass at each end, a
linear spring $EA/h$ and a Kelvin--Voigt dashpot, so that the damping matrix is
$\mat C=t_v\mat K$. The take-up adds the carriage displacement $y_c$ as a degree of
freedom, with mass $\Mtu$ and the static force $\nstr\Tt$, and enters the element that
wraps the take-up pulley through its elongation $u_R-u_L+\nstr y_c$. The weight
components of belt and material along the path and the counterweight are constant
loads, the motion resistances are applied with the same onset $\varphi$ as in the
modal solution, and the two end nodes follow the prescribed drive displacement. The
static state is solved first, and the start is integrated with the Newmark
average-acceleration scheme, which is implicit, unconditionally stable, second-order
accurate and free of numerical damping; the time step is refined in proportion to the
element length, so that the convergence of \cref{fig:validation}(d) measures the
combined error in space and time.
% src/beltloop/lumped.py; tests: test_lumped.py, test_strands.py
